\documentclass[11pt]{article}
\usepackage{mathblatt}
% Präambel des Setzers (werkzeuge/setzer-praeambel.tex, wortgleich); dazu \gteil (Muster T2U)
\makeatletter
\setlength{\mbpfbe}{0mm}% keine Punkte (pruefung.md § 5: nur beim Prüfstein)
\newlength{\gEin}\setlength{\gEin}{0pt}
\let\gSchrift\relax
\renewcommand{\pfkreuzzeile}[1]{\par\noindent\hangindent3.2em\hangafter1\hspace*{1.6em}$\square$\ #1\par}
\newcommand{\gkreuz}[1]{$\square$\ #1\hspace{2em}}
\newcommand{\gkreuzzeile}[2]{\par\noindent\hangindent2.2em\hangafter1\mbox{#1\,$\square$}\ #2\par}% Buchstabe, Option: eine je Zeile (neben Skizze)
% links, was man ansieht, rechts, was man tut; Spaltengrenze überall an derselben Stelle (gemeinsam.md „Satz“)
% v0.17: Fußeintrag „Nr: Wert“ – kurz in einer Box (nie getrennt), lang umbrechbar mit der Nummer am Anfang
% (vorher \mbox: lange Lösungen liefen über den Rand, AS40 138 pt); robust, weil \fusshilfe in die .aux schreibt
\NewDocumentCommand{\gfuss}{m m}{\sbox0{#1:~#2}\ifdim\wd0<0.45\textwidth\mbox{#1:~#2}\else\mbox{#1:}\nobreak\ #2\fi}
\newsavebox{\gbox}\newlength{\gLinks}
% v0.10 (Lehrer 13.10.): keine starre Halbierung – die linke Spalte ist so breit wie die Figur (+ 6 mm), die
% Antwort steht direkt daneben; ist die Figur breiter als 0,62 Zeile, steht die Antwort darunter
\newcommand{\gzwei}[2]{\par\smallskip\noindent\sbox{\gbox}{#1}%
  \setlength{\gLinks}{\dimexpr\wd\gbox+6mm\relax}%
  \ifdim\gLinks>0.62\linewidth
    \usebox{\gbox}\par\smallskip\noindent\begin{minipage}{\linewidth}#2\end{minipage}%
  \else
    \begin{minipage}[t]{\gLinks}\vspace{0pt}\usebox{\gbox}\end{minipage}%
    \begin{minipage}[t]{\dimexpr\linewidth-\gLinks\relax}\vspace{0pt}\raggedright #2\end{minipage}%
  \fi\par}
% v0.9 (12.10., pruefung.md § 5 „Marke mit B, Bedingungen rechts“): drittes Argument von pfaufg = Bedingungen
% der Prüfung (v0.10: \pfohnetr = TR durchgestrichen, „FOR“), klein grau im rechten Rand auf der Grundlinie der Nummer, Zeilen mit \\
\newcommand{\gbedrand}[2]{\ifblank{#2}{}{\rlap{\hspace*{\dimexpr#1+2mm\relax}%
  \parbox[t]{15mm}{\raggedright\scriptsize\leavevmode\color{mbgrau}\strut#2}}}}
% Teil eines Bündels (a, b …): Marke links im Rand, Bedingungen rechts im Rand, auf der Zeile des Teils
\newcommand{\gteilrand}[2]{\llap{\mbpfrandmarke{#1}\hspace*{\mbpfnr}}\gbedrand{\linewidth}{#2}}
% Teil eines Bündels: Buchstabe in eigener Spalte, Inhalt eingerückt (\gEin = Einzug, damit die Spaltengrenze von
% \gzwei an derselben Stelle bleibt)
\newenvironment{gbteil}[3]{\par\smallskip\noindent\gteilrand{#2}{#3}\makebox[\mbpfteil][l]{\textbf{#1}}%
  \begin{minipage}[t]{\dimexpr\linewidth-\mbpfteil\relax}\setlength{\gEin}{\mbpfteil}\raggedright\strut\ignorespaces}%
 {\end{minipage}\par}
% Grafik höchstens so breit wie die Spalte (Bündel nebeneinander, Koordinatensystem ohne Karo)
\newsavebox{\gpbox}
\newcommand{\gpassend}[1]{\sbox{\gpbox}{#1}\ifdim\wd\gpbox>\linewidth\typeout{GPASSEND: verkleinert \the\wd\gpbox\space auf \the\linewidth}\resizebox{\linewidth}{!}{\usebox{\gpbox}}\else\usebox{\gpbox}\fi}
% Skizzenfeld (v0.10, Lehrer 13.10.): Karo ab dem linken Rand über die Zeile, „Skizze“ klein grau oben links im Feld
% (wie \pfskizzenfeld der Vorlage); der rechte Rand bleibt den Bedingungen
\newcommand{\gskizze}[1]{\par\smallskip\noindent\begin{tikzpicture}
  \clip (0,0) rectangle ({\linewidth-1mm},{#1*10mm});
  \draw[black!16,line width=0.3pt,step=5mm] (0,0) grid ({\linewidth},{#1*10mm});
  \node[anchor=north west,inner sep=2pt,font=\scriptsize,text=mbgrau,fill=white] at (0,{#1*10mm}) {Skizze};
  \end{tikzpicture}\par}
\RenewDocumentEnvironment{pfaufg}{m m m}{%
  \begin{lrbox}{\mbaufgabebox}\begin{minipage}[t]{\linewidth}%
  \gSchrift\noindent\hspace*{\gEin}\mbpfrandmarke{#1}\gbedrand{\linewidth-\gEin}{#3}%
  \begin{minipage}[t]{\mbpfnr}\strut\textbf{#2}\end{minipage}%
  \begin{minipage}[t]{\dimexpr\linewidth-\gEin-\mbpfnr\relax}\raggedright\strut\ignorespaces}%
 {\end{minipage}%
  \end{minipage}\end{lrbox}%
  \setlength{\mbaufgabehoehe}{\dimexpr\ht\mbaufgabebox+\dp\mbaufgabebox\relax}%
  \ifdefempty{\gkopfnach}{\par\addvspace{5pt}\Needspace*{\mbaufgabehoehe}}{%
    % Stufenkopf, der mit dieser (ersten) Aufgabe zusammenbleibt: Platz für beide verlangen, dann den Kopf setzen
    % (12.10., Quadratische: ein Kopf stand allein unten, die Aufgabe mit Koordinatensystem auf der nächsten Seite)
    % v0.17 (AS40 „Terme deuten“ allein unten): hinter dem Kopf steht das Sprungziel (\special, ein Knoten, nach
    % dem die Seite brechen durfte) – darum \nobreak nach dem Kopf; Vorhalt 56 pt = Kopf mit Linie und Abständen
    \par\Needspace*{\dimexpr\mbaufgabehoehe+56pt\relax}\gkopfnach\par\nobreak\gdef\gkopfnach{}\addvspace{5pt}}%
  \noindent\usebox{\mbaufgabebox}\par\addvspace{5pt}}
% Stufenkopf vormerken: \gkopfvor{<Kopf>} – die nächste pfaufg setzt ihn, sobald beide zusammen Platz haben
\def\gkopfnach{}
\newcommand{\gkopfvor}[1]{\gdef\gkopfnach{#1}}
% Teilaufgabe mit Skizze links in gzwei: Buchstabe in eigener Spalte, darunter Text und Skizze
\newcommand{\gteil}[2]{\makebox[\mbpfteil][l]{\textbf{#1}}\parbox[t]{\dimexpr0.5\textwidth-\gEin-\mbpfnr-5mm-\mbpfteil\relax}{\raggedright #2}}
% Stufenkopf der Vorlage (mathblatt.sty 2026-10-12: Linie, Abstand, 13 pt, \pfjahre) ohne ihren Vorhalt von
% 14 Zeilen: der Setzer fordert vor dem Kopf 0,16 Seitenhöhe (Kopf + erste Aufgabe), mehr nicht
\setlength{\mbpfstufevorhalt}{0pt}
\makeatother
\newcommand{\karo}[1]{\par\smallskip\noindent\begin{tikzpicture}
  \clip (0,0) rectangle ({\linewidth-1mm},{#1*10mm});
  \draw[black!16,line width=0.3pt,step=5mm] (0,0) grid ({\linewidth},{#1*10mm});
  \end{tikzpicture}\par}
\newcommand{\kk}[1]{$\square$\,#1\hspace{0.9em}}
\newcommand{\linie}[1]{\,{\color{black!45}\rule[-1pt]{#1}{0.4pt}}\,}
\newcommand{\pkt}{\textperiodcentered{}}
% v0.14 (14.10.): Stufenkopf mit fünf Punkten (Option = Jahre für die Bestellkarte) oder „zuletzt 19“
\NewDocumentCommand{\gpunkte}{o m}{#2}
\newcommand{\gzuletzt}[1]{zuletzt #1}
% v0.7 (12.10.): Schülerfrage unter dem Stufenkopf; Ankreuzoption als Bild (Buchstabe und Kästchen darunter)
\newcommand{\gfrage}[1]{\par\nobreak\noindent\hspace*{\mbpfnr}{\small\itshape #1}\par\nobreak\addvspace{2pt}}
\newlength{\gbildbreite}\setlength{\gbildbreite}{0.32\linewidth}% Setzer setzt sie je Reihe (vier Bilder: 0,24)
\newcommand{\gbildoption}[3]{\begin{minipage}[t][#3][b]{\gbildbreite}\centering #2\par\smallskip\mbox{#1\,$\square$}\end{minipage}}
% Lineare (12.10.): Bild mit Buchstabe ohne Kästchen (Diagramm zuordnen); Aussage mit „wahr / falsch“ zum Ankreuzen
\newcommand{\gbildname}[3]{\begin{minipage}[t][#3][b]{\gbildbreite}\centering #2\par\smallskip\mbox{\textbf{#1}}\end{minipage}}
\newcommand{\gwahrfalsch}[2]{\par\noindent\hangindent2.2em\hangafter1\makebox[2.2em][l]{(#1)}#2\hfill\mbox{wahr $\square$\quad falsch $\square$}\par}
\begin{document}

\pfheftstil
\pflaufkopf[Abitur GK \pkt{} Prüfung]{Rechter Winkel}{AG30}\pfseitevon
\pfheftkopf{Rechter Winkel}{}
\pfstufenzeile{rechten Winkel nachweisen \pkt{} Normalenvektor nachweisen \pkt{} Parameter oder Punkt gesucht}

% ===================== Stufe: rechten Winkel nachweisen =====================
\gkopfvor{\pfstufekopf{rechten Winkel nachweisen}{\gpunkte[26 25 24]{\textcolor{black!80}{$\bullet$}\,\textcolor{black!80}{$\bullet$}\,\textcolor{black!80}{$\bullet$}\,\textcolor{black!25}{$\circ$}\,\textcolor{black!25}{$\circ$}}}\special{pdf:dest (stufe-rechten-winkel-nachweisen) [@thispage /XYZ @xpos @ypos null]}}
% kein Einstieg: keine Schülerfrage in der Gliederung (Schülerfrage, -Fälle, -Form, -Weg, -Skizze; oder Frage) – Einstieg nur von dort, nie aus der Bank (pruefung.md § 5, 12.10.)
% 2026-bb-gk-A1.5a-m1 (Original 2026-bb-gk-A1.5a, Schwere I, rang 1; herausgelöste Fassung: keine)
\begin{pfaufg}{GK~’26\\A1.5a}{1.}{\pfohnetr}
Gegeben ist das Dreieck ABC mit $A(0 \mid 0 \mid 0)$, $B(4 \mid 4 \mid 7)$ und $C(t \mid -t \mid 0)$, $t > 0$. Zeige, dass das Dreieck in $A$ einen rechten Winkel hat.
\karo{3}
\fusshilfe{\gfuss{1}{rechter Winkel in A}}
\end{pfaufg}
% 2025-bebb-gk-A1.2b-m1 (Original 2025-bebb-gk-A1.2b, Schwere I, rang 2; herausgelöste Fassung: keine)
\begin{pfaufg}{GK~’25\\A1.2b}{2.}{\pfohnetr}
Gegeben ist die Gerade $g\colon \vec{x} = \begin{pmatrix}6\\2\\-2\end{pmatrix} + s \cdot \begin{pmatrix}-3\\0\\2\end{pmatrix}$, $s \in \mathbb{R}$. Die Gerade $h$ verläuft parallel zur y-Achse und schneidet $g$ im Punkt $(6 \mid 2 \mid -2)$. Untersuche, ob $g$ und $h$ senkrecht zueinander verlaufen.
\karo{4}
\fusshilfe{\gfuss{2}{ja}}
\end{pfaufg}
% 2024-bebb-gk-B3b-m1 (Original 2024-bebb-gk-B3b, Schwere I, rang 3; herausgelöste Fassung: keine)
\begin{pfaufg}{GK~’24\\B3b}{3.}{}
Gegeben ist die Pyramide ABCDS mit $A(0 \mid 0 \mid 0)$, $B(3 \mid 0 \mid 0)$, $C(3 \mid 3 \mid 0)$, $D(0 \mid 6 \mid 0)$ und $S(0 \mid 0 \mid 5)$. Zeige, dass das Dreieck CDS rechtwinklig ist.
\gzwei{\begin{ksys3}[xyz,x1min=0,x1max=3,x2min=0,x2max=6,x3min=0,x3max=5,karo=0.22]\rstrecke[dashed]{0,0,0}{3,0,0}\rstrecke{3,0,0}{3,3,0}\rstrecke{3,3,0}{0,6,0}\rstrecke[dashed]{0,6,0}{0,0,0}\rstrecke[dashed]{0,0,0}{0,0,5}\rstrecke{3,0,0}{0,0,5}\rstrecke{3,3,0}{0,0,5}\rstrecke{0,6,0}{0,0,5}\rpunkt*[above left]{0}{0}{0}{A}\rpunkt*[below left]{3}{0}{0}{B}\rpunkt*[below right]{3}{3}{0}{C}\rpunkt*[above right]{0}{6}{0}{D}\rpunkt*[above left]{0}{0}{5}{S}\end{ksys3}}{\karo{4}}
\fusshilfe{\gfuss{3}{rechter Winkel bei C}}
\end{pfaufg}

% ===================== Stufe: Normalenvektor nachweisen =====================
\gkopfvor{\pfstufekopf{Normalenvektor nachweisen}{\gpunkte[25 22]{\textcolor{black!25}{$\circ$}\,\textcolor{black!80}{$\bullet$}\,\textcolor{black!25}{$\circ$}\,\textcolor{black!25}{$\circ$}\,\textcolor{black!80}{$\bullet$}}}\special{pdf:dest (stufe-normalenvektor-nachweisen) [@thispage /XYZ @xpos @ypos null]}}
% kein Einstieg: keine Schülerfrage in der Gliederung (Schülerfrage, -Fälle, -Form, -Weg, -Skizze; oder Frage) – Einstieg nur von dort, nie aus der Bank (pruefung.md § 5, 12.10.)
% 2025-bebb-gk-A1.8a-m1 (Original 2025-bebb-gk-A1.8a, Schwere I, rang 4; herausgelöste Fassung: keine)
\begin{pfaufg}{GK~’25\\A1.8a}{4.}{\pfohnetr}
Gegeben ist die Ebene $E\colon \vec{x} = \begin{pmatrix}2\\-1\\0\end{pmatrix} + r \cdot \begin{pmatrix}-4\\3\\1\end{pmatrix} + s \cdot \begin{pmatrix}4\\-3\\0\end{pmatrix}$, $r, s \in \mathbb{R}$. Zeige, dass der Vektor $\left(\begin{smallmatrix}3\\4\\0\end{smallmatrix}\right)$ senkrecht auf $E$ steht.
\karo{4}
\fusshilfe{\gfuss{4}{senkrecht auf E}}
\end{pfaufg}
% 2022-bebb-gk-B3c-m1 (Original 2022-bebb-gk-B3c, Schwere II, rang 5; herausgelöste Fassung: keine)
\begin{pfaufg}{GK~’22\\B3c}{5.}{}
Ein Kirchturmdach steht auf dem Quadrat ABCD mit $A(0 \mid 0 \mid 0)$, $B(6 \mid 0 \mid 0)$, $C(6 \mid 6 \mid 0)$, $D(0 \mid 6 \mid 0)$. Darüber liegen $E(3 \mid 0 \mid 5)$, $F(6 \mid 3 \mid 5)$, $G(3 \mid 6 \mid 5)$, $H(0 \mid 3 \mid 5)$ und die Spitze $S(3 \mid 3 \mid 10)$; 1 LE entspricht 1 m. Das Dach hat vier gleiche viereckige Dachflächen wie CGSF und vier dreieckige Giebelflächen. Die Ebene L enthält C, G und F; die Ebene N ist parallel zu L und enthält B. Weise nach, dass $\left(\begin{smallmatrix}5\\5\\3\end{smallmatrix}\right)$ ein Normalenvektor von N ist. Bestimme eine Koordinatengleichung von N.
\karo{4}
\fusshilfe{\gfuss{5}{Normalenvektor von L und N; 30}}
\end{pfaufg}

% ===================== Stufe: Parameter oder Punkt gesucht =====================
\gkopfvor{\pfstufekopf{Parameter oder Punkt gesucht}{\gpunkte[26 25 24 22]{\textcolor{black!80}{$\bullet$}\,\textcolor{black!80}{$\bullet$}\,\textcolor{black!80}{$\bullet$}\,\textcolor{black!25}{$\circ$}\,\textcolor{black!80}{$\bullet$}}}\special{pdf:dest (stufe-parameter-oder-punkt-gesucht) [@thispage /XYZ @xpos @ypos null]}}
% kein Einstieg: keine Schülerfrage in der Gliederung (Schülerfrage, -Fälle, -Form, -Weg, -Skizze; oder Frage) – Einstieg nur von dort, nie aus der Bank (pruefung.md § 5, 12.10.)
% 2022-bebb-gk-A1.4c-m1 (Original 2022-bebb-gk-A1.4c, Schwere I, rang 6; herausgelöste Fassung: keine)
\begin{pfaufg}{GK~’22\\A1.4c}{6.}{\pfohnetr}
Gegeben sind die Ebenen $E\colon 2x - 5y = 0$ und $F\colon 3x + sy + z = 4$. Bestimme $s$ so, dass $F$ senkrecht auf $E$ steht.
\karo{3}
\fusshilfe{\gfuss{6}{1,2}}
\end{pfaufg}
% 2024-bebb-gk-A1.2b-m1 (Original 2024-bebb-gk-A1.2b, Schwere II, rang 7; herausgelöste Fassung: keine)
\begin{pfaufg}{GK~’24\\A1.2b}{7.}{\pfohnetr}
Der Ursprung O, $P(1 \mid 0 \mid 26)$ und $Q(5 \mid t \mid 23)$ bilden ein Dreieck. Ermittle die Werte von $t$, für die das Dreieck in $Q$ einen rechten Winkel hat.
\karo{3}
\fusshilfe{\gfuss{7}{±7}}
\end{pfaufg}
% 2025-bebb-gk-B3d-m1 (Original 2025-bebb-gk-B3d, Schwere II, rang 8; herausgelöste Fassung: keine)
\begin{pfaufg}{GK~’25\\B3d}{8.}{}
Gegeben ist die Pyramide ABCDS mit $A(0 \mid 0 \mid 0)$, $B(3 \mid 3 \mid 0)$, $C(0 \mid 8 \mid 0)$, $D(-3 \mid 3 \mid 0)$ und $S(0 \mid 0 \mid 6)$; ihre Grundfläche ABCD ist ein Drachenviereck. Der Punkt B wird parallel zur y-Achse nach B' verschoben, so dass das Viereck AB'CD in B' einen rechten Innenwinkel hat. Erläutere den Ansatz $\left(\begin{smallmatrix}3\\y\\0\end{smallmatrix}\right)$ · $\left(\begin{smallmatrix}3\\y  -  8\\0\end{smallmatrix}\right)$ = 0.
\gzwei{\begin{ksys3}[xyz,x1min=-3,x1max=3,x2min=0,x2max=8,x3min=0,x3max=6,karo=0.2]\rstrecke{0,0,0}{3,3,0}\rstrecke{3,3,0}{0,8,0}\rstrecke[dashed]{0,8,0}{-3,3,0}\rstrecke[dashed]{-3,3,0}{0,0,0}\rstrecke{0,0,0}{0,0,6}\rstrecke{3,3,0}{0,0,6}\rstrecke{0,8,0}{0,0,6}\rstrecke[dashed]{-3,3,0}{0,0,6}\rpunkt*[below left]{0}{0}{0}{A}\rpunkt*[below]{3}{3}{0}{B}\rpunkt*[below right]{0}{8}{0}{C}\rpunkt*[above]{-3}{3}{0}{D}\rpunkt*[left]{0}{0}{6}{S}\end{ksys3}}{\karo{4}}
\fusshilfe{\gfuss{8}{rechter Winkel bei B'}}
\end{pfaufg}
% 2026-bb-gk-A1.8a-m1 (Original 2026-bb-gk-A1.8a, Schwere III, rang 9; herausgelöste Fassung: keine)
\begin{pfaufg}{GK~’26\\A1.8a}{9.}{\pfohnetr}
Gegeben sind $A(1 \mid -2 \mid 0)$ und $B(7 \mid 1 \mid -3)$. Für den Punkt $C$ gilt: Seine z-Koordinate ist 5; die Geraden AB und AC stehen senkrecht zueinander; $C$ liegt in der Ebene $x + 2y = 0$. Ermittle die Koordinaten von $C$.
\karo{4}
\fusshilfe{\gfuss{9}{C($\tfrac{10}{3}$|$-\tfrac{5}{3}$|5)}}
\end{pfaufg}

\pfblattende{10}{26 & 25 & 24 & 23 & 22 & 21 & 20 & 19 & 18 & 17}{A1.5a A1.8a & A1.2b A1.8a B3d & A1.2b B3b &  & A1.4c B3c & B3a & B3.1f &  & B2.1b & }%
{\pfausschnitt{}{Rechter Winkel}{Winkel \pkt{} Abstände}}
\end{document}